% sec4_2.tex — Section 4.2: Multiple-Equation GMM Defined

\section{Multiple-Equation GMM Defined}
\label{sec:4.2}

Derivation of the GMM estimator for multiple equations, too, can be carried out in
much the same way as in the single-equation case. Let $\tilde{\boldsymbol{\delta}}$
be a hypothetical value of the true parameter vector $\boldsymbol{\delta}$ and define
$\mathbf{g}_n(\tilde{\boldsymbol{\delta}})$ by~\eqref{eq:4.1.1}. The definition of the GMM
estimator is the same as in (3.4.5), provided that the weighting matrix
$\hat{\mathbf{W}}$ is now
$\sum_m K_m \times \sum_m K_m$. In the previous section, we were able to rewrite
$\E[\mathbf{g}(\mathbf{w}_i;\, \tilde{\boldsymbol{\delta}})]$ for multiple equations as a
linear function of $\tilde{\boldsymbol{\delta}}$ (see~\eqref{eq:4.1.8}).
We can do the same for the sample analogue
$\mathbf{g}_n(\tilde{\boldsymbol{\delta}})$:
\begin{align}
\underset{(\sum_{m=1}^{M} K_m \times 1)}{\mathbf{g}_n(\tilde{\boldsymbol{\delta}})}
&= \begin{bmatrix}
\dfrac{1}{n} \displaystyle\sum_{i=1}^{n} \mathbf{x}_{i1} \cdot
(y_{i1} - \mathbf{z}_{i1}' \tilde{\boldsymbol{\delta}}_1) \\[6pt]
\vdots \\[4pt]
\dfrac{1}{n} \displaystyle\sum_{i=1}^{n} \mathbf{x}_{iM} \cdot
(y_{iM} - \mathbf{z}_{iM}' \tilde{\boldsymbol{\delta}}_M)
\end{bmatrix} \notag \\[6pt]
&= \begin{bmatrix}
\dfrac{1}{n} \displaystyle\sum_{i=1}^{n} \mathbf{x}_{i1} \cdot y_{i1} \\[6pt]
\vdots \\[4pt]
\dfrac{1}{n} \displaystyle\sum_{i=1}^{n} \mathbf{x}_{iM} \cdot y_{iM}
\end{bmatrix}
- \begin{bmatrix}
\dfrac{1}{n} \displaystyle\sum_{i=1}^{n}
\mathbf{x}_{i1} \mathbf{z}_{i1}' \tilde{\boldsymbol{\delta}}_1 \\[6pt]
\vdots \\[4pt]
\dfrac{1}{n} \displaystyle\sum_{i=1}^{n}
\mathbf{x}_{iM} \mathbf{z}_{iM}' \tilde{\boldsymbol{\delta}}_M
\end{bmatrix} \notag \\[6pt]
&= \begin{bmatrix}
\dfrac{1}{n} \displaystyle\sum_{i=1}^{n} \mathbf{x}_{i1} \cdot y_{i1} \\[6pt]
\vdots \\[4pt]
\dfrac{1}{n} \displaystyle\sum_{i=1}^{n} \mathbf{x}_{iM} \cdot y_{iM}
\end{bmatrix}
- \begin{bmatrix}
\dfrac{1}{n} \displaystyle\sum_{i=1}^{n} \mathbf{x}_{i1} \mathbf{z}_{i1}' & & \\
& \ddots & \\
& & \dfrac{1}{n} \displaystyle\sum_{i=1}^{n} \mathbf{x}_{iM} \mathbf{z}_{iM}'
\end{bmatrix}
\begin{bmatrix}
\tilde{\boldsymbol{\delta}}_1 \\ \vdots \\ \tilde{\boldsymbol{\delta}}_M
\end{bmatrix} \notag \\[6pt]
&\equiv \mathbf{s}_{\mathbf{x}\mathbf{y}}
- \mathbf{S}_{\mathbf{x}\mathbf{z}}\, \tilde{\boldsymbol{\delta}},
\label{eq:4.2.1}
\end{align}
where
\begin{equation}
\mathbf{s}_{\mathbf{x}\mathbf{y}}
\equiv \begin{bmatrix}
\dfrac{1}{n} \displaystyle\sum_{i=1}^{n} \mathbf{x}_{i1} \cdot y_{i1} \\[6pt]
\vdots \\[4pt]
\dfrac{1}{n} \displaystyle\sum_{i=1}^{n} \mathbf{x}_{iM} \cdot y_{iM}
\end{bmatrix}, \quad
\mathbf{S}_{\mathbf{x}\mathbf{z}}
\equiv \begin{bmatrix}
\dfrac{1}{n} \displaystyle\sum_{i=1}^{n} \mathbf{x}_{i1} \mathbf{z}_{i1}' & & \\
& \ddots & \\
& & \dfrac{1}{n} \displaystyle\sum_{i=1}^{n} \mathbf{x}_{iM} \mathbf{z}_{iM}'
\end{bmatrix}.
\label{eq:4.2.2}
\end{equation}

Again, this is the \textit{same} in form as the expression for
$\mathbf{g}_n(\tilde{\boldsymbol{\delta}})$ for the single-equation case.

It follows that the results of Section~3.4 about the derivation of the GMM estimator
are applicable to the multiple-equation case. In particular, just reproducing
(3.4.8) and (3.4.11),
\begin{align}
\text{multiple-equation GMM estimator:}\quad
\hat{\boldsymbol{\delta}}(\hat{\mathbf{W}})
&= (\mathbf{S}_{\mathbf{x}\mathbf{z}}' \hat{\mathbf{W}} \mathbf{S}_{\mathbf{x}\mathbf{z}})^{-1}
\mathbf{S}_{\mathbf{x}\mathbf{z}}' \hat{\mathbf{W}} \mathbf{s}_{\mathbf{x}\mathbf{y}},
\label{eq:4.2.3} \\[4pt]
\text{its sampling error:}\quad
\hat{\boldsymbol{\delta}}(\hat{\mathbf{W}}) - \boldsymbol{\delta}
&= (\mathbf{S}_{\mathbf{x}\mathbf{z}}' \hat{\mathbf{W}} \mathbf{S}_{\mathbf{x}\mathbf{z}})^{-1}
\mathbf{S}_{\mathbf{x}\mathbf{z}}' \hat{\mathbf{W}} \bar{\mathbf{g}}.
\label{eq:4.2.4}
\end{align}

The features specific to the multiple-equation GMM formula are the following:
(i)~$\mathbf{s}_{\mathbf{x}\mathbf{y}}$ is a stacked vector as in~\eqref{eq:4.2.2},
(ii)~$\mathbf{S}_{\mathbf{x}\mathbf{z}}$ is a block diagonal matrix as in~\eqref{eq:4.2.2},
(iii)~accordingly, the size of the weighting matrix $\hat{\mathbf{W}}$ is
$\sum_m K_m \times \sum_m K_m$,
and (iv)~$\bar{\mathbf{g}}$ in~\eqref{eq:4.2.4}, the sample mean of $\mathbf{g}_i$,
is the stacked vector
\begin{equation}
\bar{\mathbf{g}} \equiv \frac{1}{n} \sum_{i=1}^{n} \mathbf{g}_i
= \begin{bmatrix}
\dfrac{1}{n} \displaystyle\sum_{i=1}^{n} \mathbf{x}_{i1} \cdot \varepsilon_{i1} \\[6pt]
\vdots \\[4pt]
\dfrac{1}{n} \displaystyle\sum_{i=1}^{n} \mathbf{x}_{iM} \cdot \varepsilon_{iM}
\end{bmatrix}
= \mathbf{g}_n(\boldsymbol{\delta}).
\label{eq:4.2.5}
\end{equation}

It will be necessary to rewrite the multiple-equation GMM formula~\eqref{eq:4.2.3}
while explicitly taking account of those special features. If
$\hat{\mathbf{W}}_{mh}$ $(K_m \times K_h)$ is the $(m, h)$ block of
$\hat{\mathbf{W}}$ $(m, h = 1, 2, \ldots, M)$, then~\eqref{eq:4.2.3} can be written
out in full as
\begin{align}
\hat{\boldsymbol{\delta}}(\hat{\mathbf{W}})
&= \begin{bmatrix}
\hat{\boldsymbol{\delta}}_1(\hat{\mathbf{W}}) \\ \vdots \\
\hat{\boldsymbol{\delta}}_M(\hat{\mathbf{W}})
\end{bmatrix} \notag \\[6pt]
&= \left(
\begin{bmatrix}
\dfrac{1}{n} \displaystyle\sum \mathbf{z}_{i1} \mathbf{x}_{i1}' & & \\
& \ddots & \\
& & \dfrac{1}{n} \displaystyle\sum \mathbf{z}_{iM} \mathbf{x}_{iM}'
\end{bmatrix}
\begin{bmatrix}
\hat{\mathbf{W}}_{11} & \cdots & \hat{\mathbf{W}}_{1M} \\
\vdots & & \vdots \\
\hat{\mathbf{W}}_{M1} & \cdots & \hat{\mathbf{W}}_{MM}
\end{bmatrix}
\begin{bmatrix}
\dfrac{1}{n} \displaystyle\sum \mathbf{x}_{i1} \mathbf{z}_{i1}' & & \\
& \ddots & \\
& & \dfrac{1}{n} \displaystyle\sum \mathbf{x}_{iM} \mathbf{z}_{iM}'
\end{bmatrix}
\right)^{-1} \notag \\[6pt]
&\quad \times
\begin{bmatrix}
\dfrac{1}{n} \displaystyle\sum \mathbf{z}_{i1} \mathbf{x}_{i1}' & & \\
& \ddots & \\
& & \dfrac{1}{n} \displaystyle\sum \mathbf{z}_{iM} \mathbf{x}_{iM}'
\end{bmatrix}
\begin{bmatrix}
\hat{\mathbf{W}}_{11} & \cdots & \hat{\mathbf{W}}_{1M} \\
\vdots & & \vdots \\
\hat{\mathbf{W}}_{M1} & \cdots & \hat{\mathbf{W}}_{MM}
\end{bmatrix}
\begin{bmatrix}
\dfrac{1}{n} \displaystyle\sum \mathbf{x}_{i1} \cdot y_{i1} \\[6pt]
\vdots \\[4pt]
\dfrac{1}{n} \displaystyle\sum \mathbf{x}_{iM} \cdot y_{iM}
\end{bmatrix} \notag \\[6pt]
&= \begin{bmatrix}
\Bigl(\dfrac{1}{n} \displaystyle\sum \mathbf{z}_{i1} \mathbf{x}_{i1}'\Bigr)
\hat{\mathbf{W}}_{11}
\Bigl(\dfrac{1}{n} \displaystyle\sum \mathbf{x}_{i1} \mathbf{z}_{i1}'\Bigr)
& \cdots &
\Bigl(\dfrac{1}{n} \displaystyle\sum \mathbf{z}_{i1} \mathbf{x}_{i1}'\Bigr)
\hat{\mathbf{W}}_{1M}
\Bigl(\dfrac{1}{n} \displaystyle\sum \mathbf{x}_{iM} \mathbf{z}_{iM}'\Bigr) \\[6pt]
\vdots & & \vdots \\[4pt]
\Bigl(\dfrac{1}{n} \displaystyle\sum \mathbf{z}_{iM} \mathbf{x}_{iM}'\Bigr)
\hat{\mathbf{W}}_{M1}
\Bigl(\dfrac{1}{n} \displaystyle\sum \mathbf{x}_{i1} \mathbf{z}_{i1}'\Bigr)
& \cdots &
\Bigl(\dfrac{1}{n} \displaystyle\sum \mathbf{z}_{iM} \mathbf{x}_{iM}'\Bigr)
\hat{\mathbf{W}}_{MM}
\Bigl(\dfrac{1}{n} \displaystyle\sum \mathbf{x}_{iM} \mathbf{z}_{iM}'\Bigr)
\end{bmatrix}^{-1} \notag \\[6pt]
&\quad \times
\begin{bmatrix}
\Bigl(\dfrac{1}{n} \displaystyle\sum \mathbf{z}_{i1} \mathbf{x}_{i1}'\Bigr)
\hat{\mathbf{W}}_{11}
\Bigl(\dfrac{1}{n} \displaystyle\sum \mathbf{x}_{i1} y_{i1}\Bigr)
+ \cdots +
\Bigl(\dfrac{1}{n} \displaystyle\sum \mathbf{z}_{i1} \mathbf{x}_{i1}'\Bigr)
\hat{\mathbf{W}}_{1M}
\Bigl(\dfrac{1}{n} \displaystyle\sum \mathbf{x}_{iM} y_{iM}\Bigr) \\[6pt]
\vdots \\[4pt]
\Bigl(\dfrac{1}{n} \displaystyle\sum \mathbf{z}_{iM} \mathbf{x}_{iM}'\Bigr)
\hat{\mathbf{W}}_{M1}
\Bigl(\dfrac{1}{n} \displaystyle\sum \mathbf{x}_{i1} y_{i1}\Bigr)
+ \cdots +
\Bigl(\dfrac{1}{n} \displaystyle\sum \mathbf{z}_{iM} \mathbf{x}_{iM}'\Bigr)
\hat{\mathbf{W}}_{MM}
\Bigl(\dfrac{1}{n} \displaystyle\sum \mathbf{x}_{iM} y_{iM}\Bigr)
\end{bmatrix}.
\label{eq:4.2.6}
\end{align}
To go from the second-to-last to the last line, use formulas (A.6) and (A.7) of
the appendix on partitioned matrices. If you find the operation too much, just set
$M = 2$.
